\chapter{Kinds of Comments}

%------------------------------------------------------------------------------

\section{The Syntax}
Syntactically, the \Java programming language knows the following types of comments:
\begin{itemize}[nosep]
	\item{\nameref{sec:BlockComments}\index{block comment}}
	\item{\nameref{sec:SingleLineComments}\index{single-line comment}}
	\item{\hyperref[sec:TrailingComments]{Trailing or end-of-line Comments}\index{trailing comment}}
\end{itemize}

%------------------------------------------------------------------------------

\subsection{Block Comments}\label{sec:BlockComments}
A \textit{block comment}\index{block comment} is delimited by \bdq{\lstinline|/*|} and \bdq{\lstinline|*/|}, everything inbetween belongs to the comment. The indentation of the opening tag is aligned with the indentation of the surrounding code, and if the block comment spans multiple lines, the asterisk of the closing tag is aligned with the asterisk of the opening tag.

For a multi-line block comment\index{block comment}, no text will be placed to the lines with the opening and the closing tag, and usually, each comment line starts with an asterisk that is aligned with that of the opening tag.

\keeplisting{3}
This looks like shown below:
\begin{lstlisting}
/* A block-comment with just one line */

/*
 * A multi-line block comment.
 */
\end{lstlisting}

Instead of a block comment\index{block comment} with just one line, you should preferably use a single-line comment\index{single-line comment}. But one-line block comments comes often handy when you have to place a comment into an otherwise empty code block:
\begin{lstlisting}
public MyClass() { /* Constructor just exists */ }
\end{lstlisting}

A block comment also can start in the same line after a statement, but that is not recommended. Nevertheless, it is sometimes handy in cases when a trailing or end-of-line comment\index{trailing comment} gets too long; see \hyperref[sec:TrailingComments]{below} for a sample.

%------------------------------------------------------------------------------

\subsection{Single Line Comments}\label{sec:SingleLineComments}
In Java, a \textit{single-line comment}\index{single-line comment} begins with \bdq{\lstinline|//|} and ends at the end of the line. This means it can be placed after a statement, in the same line; in that case, we talk about a \hyperref[sec:TrailingComments]{trailing or end-of-line comment}\index{trailing comment}, as described \hyperref[sec:TrailingComments]{below}.

\keeplisting{7}
I use single-line comments\index{single-line comment} mainly as headings, so I gave them a special format:
\begin{lstlisting}
//---* A heading *---------------------------------------------------
…

//---* Done! *-------------------------------------------------------
return;
\end{lstlisting}

Longer comments should be placed always into a block comment\index{block comment} instead of spanning it over multiple single-line comments\index{single-line comment}. 

%------------------------------------------------------------------------------

\subsection{Trailing Comments}\label{sec:TrailingComments}
\keeplisting{2}
A \textit{trailing} or \textit{end-of-line comment}\index{trailing comment} is technically a single-line comment\index{single-line comment} that is placed after a statement in the same line:
\begin{lstlisting}
final var localVariable; // Short description of localVariable
\end{lstlisting}

When the comment gets too long (exceeding column~80) it should be continued in the following line:
\begin{lstlisting}
final var localVariable; // This is a much, much longer description of 
    // the localVariable
\end{lstlisting}

In such a case, you could consider also to use a block comment\index{block comment} instead, like this:
\begin{lstlisting}
final var localVariable; /* This is a much, much longer description of 
    the localVariable */
\end{lstlisting}

Another consideration is whether the comment should be a trailing comment\index{trailing comment} at all, although this depends from the context.

%------------------------------------------------------------------------------

\section{The Semantics}
Semantically, Java source code can have generally four kinds of comments:
\begin{itemize}[nosep]
	\item{\nameref{sec:StructuringComments}}
	\item{\nameref{sec:DocumentationComments}}
	\item{\nameref{sec:ImplementationComments}}
	\item{\nameref{sec:MaintenanceComments}}
\end{itemize}

The \textit{structuring comments} were mentioned already in chapter \tqfullvref{sec:ClassAndInterfaceDeclarations} when the skeletons for the various kinds of compilation units\index{compilation unit} where discussed.

In Java, the \textit{documentation comments} (also known as \bdq{doc comments} or \bdq{the Javadoc}) will be externalised for the generation of the program/library documentation by the Javadoc\index{javadoc} tool; this allows to create some of the program documentation together with the source code and to keep both closely together. 

Obviously, \textit{implementation comments} are the other comments, that are not externalized and published. Roughly, the documentation comments describe how to use the code (the API), unrelated to the implementation, while the implementation comments describe what the code is doing and why.

\textit{Maintenance comments} are technically a special form of implementation comments, but as they have a special function, they are covered separately in chapter \tqfullvref{sec:MaintenanceComments}.

All are discussed in more details in the following chapters.
  
%==============================================================================
% End of file!!
%==============================================================================

